% !TEX encoding = UTF-8 Unicode
\documentclass[10pt,aspectratio=169,]{beamer}
\setbeamercovered{transparent=10}
\usetheme[
%  showheader,
%  red,
  purple,
%  gray,
%  graytitle,
  colorblocks,
%  noframetitlerule,
]{Verona}

\usepackage[T1]{fontenc}
\usepackage[utf8]{inputenc}
\usepackage{lipsum}
%%%%%%%%%%%%%%%%%%%%%%%%%%%%%%%
% Mac上使用如下命令声明隶书字体，windows也有相关方式，大家可自行修改
%\providecommand{\lishu}{\CJKfamily{zhli}}
%%%%%%%%%%%%%%%%%%%%%%%%%%%%%%%
\usepackage{multicol}
\usepackage{tikz}
\usetikzlibrary{fadings}
%
%\setbeamertemplate{sections/subsections in toc}[ball]
\usepackage{xeCJK}
\usepackage{ctex, hyperref}
\usepackage{fontspec} 
\usepackage{listings}
\usepackage{caption}
\usepackage{subcaption}
\usefonttheme{professionalfonts}
\def\mathfamilydefault{\rmdefault}
\usepackage{amsmath}
\usepackage{multirow}
\usepackage{booktabs}
\usepackage{bm}
\usepackage[ruled,linesnumbered]{algorithm2e}
\usepackage{listings} 
\setbeamertemplate{section in toc}{\hspace*{1em}\inserttocsectionnumber.~\inserttocsection\par}
\setbeamertemplate{subsection in toc}{\hspace*{2em}\inserttocsectionnumber.\inserttocsubsectionnumber.~\inserttocsubsection\par}
\setbeamerfont{subsection in toc}{size=\small}
\AtBeginSection[]{
	\begin{frame}%
        \begin{multicols}{2}%分列两栏的目录
		    \frametitle{Outline}%
		    \kaishu{\textbf{\tableofcontents[currentsection]}}
        \end{multicols}
	\end{frame}
}
%\AtBeginSubsection[]{%
%	\begin{frame}%
%		\frametitle{Outline}%
%		\textbf{\tableofcontents[currentsection, currentsubsection]} %
%	\end{frame}%
%}
\title{大标题}
\subtitle{副标题}
\author[Carpenter Cheung]{姓名}
\mail{\url{邮箱地址}}
\institute[Nanjing normal university]{指导老师: xxx老师\\强化培养学院\\南京师范大学}
\date{\today}
\titlegraphic[width=3cm]{nnu_big_logo}{}


% 配置 listings 包的一些基本样式
\lstset{
	basicstyle=\ttfamily\small,
	keywordstyle=\color{blue},
	stringstyle=\color{red},
	commentstyle=\color{green},
	morecomment=[l][\color{magenta}]{\#},
	frame=single,
	breaklines=true,
  postbreak=\mbox{\textcolor{red}{$\hookrightarrow$}\space},
}
%%%%%%%%%%%%%%%%%%%%%%%%%%%%%%%%
% ----------- 标题页 ------------
%%%%%%%%%%%%%%%%%%%%%%%%%%%%%%%%



\begin{document}

\maketitle

%%% define code
%\defverbatim[colored]\lstI{
%	\begin{lstlisting}[language=C++,basicstyle=\ttfamily,keywordstyle=\color{red}]
%	int main() {
%	// Define variables at the beginning
%	// of the block, as in C:
%	CStash intStash, stringStash;
%	int i;
%	char* cp;
%	ifstream in;
%	string line;
%	[...]
%	\end{lstlisting}
%}
%%%%%%%%%%%%%%%%%%%%%%%%%%%%%%%%
% ----------- FRAME ------------
%%%%%%%%%%%%%%%%%%%%%%%%%%%%%%%%

\section{简介}
\subsection{研究背景}
\begin{frame}[c]{简介}
	\framesubtitle{研究背景}
\begin{block}{研究背景及场景}
\begin{itemize}
	\item ABC
\end{itemize}
\end{block}

\begin{block}{研究现状}
	\begin{itemize}
	\item 123456
	\end{itemize}
\end{block}
\songti{如果你并不是很中意阴影, 圆角, 颜色等设置, try to change them!}
\end{frame}	

%\begin{exampleblock}{Example Block}
%	Content of an example block
%\end{exampleblock}
%
%\begin{alertblock}{Alert block}
%	Content of an alert block
%\end{alertblock}

\subsection{研究目标}
\begin{frame}[c]{简介}
	\framesubtitle{研究目标}
	\begin{block}{研究目标}
		\begin{enumerate}
			\item enumerate是圆形的.
		\end{enumerate}
	\end{block}
	
	\begin{alertblock}{挑战}
		\begin{itemize}
            \item itemize是三角形.
        \end{itemize}
	\end{alertblock}
	%\begin{exampleblock}{Example Block}
	%	Content of an example block
	%\end{exampleblock}
	%
	%\begin{alertblock}{Alert block}
	%	Content of an alert block
	%\end{alertblock}
	
\end{frame}	

\section{研究方法}

\begin{frame}{研究方法}
	\framesubtitle{模型}
	\begin{equation*}
		J(\theta) = \mathbb{E}_{\pi_\theta}[G_t] = \sum_{s\in\mathcal{S}} d^\pi (s)V^\pi(s)=\sum_{s\in\mathcal{S}} d^\pi(s)\sum_{a\in\mathcal{A}}\pi_\theta(a|s)Q^\pi(s,a)
	\end{equation*}
	
	\begin{align}
		Q_\mathrm{target}&=r+\gamma Q^\pi(s^\prime, \pi_\theta(s^\prime)+\epsilon)\\
		\epsilon&\sim\mathrm{clip}(\mathcal{N}(0, \sigma), -c, c)\nonumber
	\end{align}
\end{frame}

\begin{frame}{模型}
	\begin{alertblock}{编号多行公式}
		% Taken from Mathmode.tex
		\begin{multline}
			A=\lim_{n\rightarrow\infty}\Delta x\left(a^{2}+\left(a^{2}+2a\Delta x+\left(\Delta x\right)^{2}\right)\right.\label{eq:reset}\\
			+\left(a^{2}+2\cdot2a\Delta x+2^{2}\left(\Delta x\right)^{2}\right)\\
			+\left(a^{2}+2\cdot3a\Delta x+3^{2}\left(\Delta x\right)^{2}\right)\\
			+\ldots\\
			\left.+\left(a^{2}+2\cdot(n-1)a\Delta x+(n-1)^{2}\left(\Delta x\right)^{2}\right)\right)\\
			=\frac{1}{3}\left(b^{3}-a^{3}\right)
		\end{multline}
	\end{alertblock}
\end{frame}

\begin{frame}{算法(pseudocode)}
\begin{algorithm}[H]
    \caption{Simulation-optimization heuristic}\label{algorithm}
    \KwData{current period $t$, initial inventory $I_{t-1}$, initial capital $B_{t-1}$, demand samples}
    \KwResult{Optimal order quantity $Q^{\ast}_{t}$}
    $r\leftarrow t$\;
    $\Delta B^{\ast}\leftarrow -\infty$\;
    \While{$\Delta B\leq \Delta B^{\ast}$ and $r\leq T$}{$Q\leftarrow\arg\max_{Q\geq 0}\Delta B^{Q}_{t,r}(I_{t-1},B_{t-1})$\;
    $\Delta B\leftarrow \Delta B^{Q}_{t,r}(I_{t-1},B_{t-1})/(r-t+1)$\;
    \If{$\Delta B\geq \Delta B^{\ast}$}{$Q^{\ast}\leftarrow Q$\;
    $\Delta B^{\ast}\leftarrow \Delta B$\;}
    $r\leftarrow r+1$\;}
\end{algorithm}
\end{frame}

\begin{frame}{算法(pseudocode)}
    \begin{itemize}[<+->]
        \item \kaishu{关于使用algorithm2e还是algpseudocode, 仁者乐山, 智者乐水.}
        \item \songti{本文选用的是algorithm2e.}
        \item \fangsong{好处是灵活性足够好, 很紧凑, 有线条.}
        \item \heiti{坏处是太乱了!!}
        \item \kaishu{\textbf{\color{red} 关于如何插入如matlab, python, cpp等代码, 请自行google之. 已添加listings包并进行了初步设定.}}
    \end{itemize}
\end{frame}
    
\begin{frame}{字体选择}
    \kaishu{这是楷体. }\pause

    \songti{这是宋体. }\pause

    \heiti{这是黑体. }\pause

    \fangsong{这是仿宋. }\pause

    \songti{这些是自带的中文字体, 若想用其余字体, 请自行搜寻. 事实上, 这是trivial的. }
\end{frame}

\section{应用与结果分析}
\begin{frame}{模型}
	\begin{itemize}
		\item \songti{\LaTeX 广泛用于学术界，期刊会议论文模板}
	\end{itemize}
	%% 表格演示
	\begin{table}[h]
		\centering
		\begin{tabular}{c|c}
			Microsoft\textsuperscript{\textregistered}  Word & \LaTeX \\
			\hline
			\kaishu{文字处理工具} & \fangsong{专业排版软件} \\
			\kaishu{容易上手，简单直观} & \fangsong{容易上手} \\
			\kaishu{所见即所得} & \fangsong{所见即所想，所想即所得} \\
			\kaishu{高级功能不易掌握} & \fangsong{进阶难，但一般用不到} \\
			\kaishu{处理长文档需要丰富经验} & \fangsong{和短文档处理基本无异} \\
			\kaishu{花费大量时间调格式} & \fangsong{无需担心格式，专心作者内容} \\
			\kaishu{公式排版差强人意} & \fangsong{尤其擅长公式排版} \\
			\kaishu{二进制格式，兼容性差} & \fangsong{文本文件，易读、稳定} \\
			\kaishu{付费商业许可} & \fangsong{自由免费使用} \\
		\end{tabular}
	\end{table}

\end{frame}

\section{总结与期望}
\begin{frame}{总结与期望}
	\begin{figure}[htbp]
		\centering
		\includegraphics[width=0.3\linewidth]{nnu_big_logo.jpeg}
		\caption{图片实例}
		\label{fig:1}
	\end{figure}
    
    \centering \songti{像素太低了!!!!!}
\end{frame}

{
    \setbeamercolor{background canvas}{bg=nnucolor}
\begin{frame}[plain]
	\vfill
	\centering
	{
		\centering \color{white} \Large {\textbf{\kaishu {感谢各位老师的倾听!\\恳请各位老师批评指正!}}}
        
	}
	\vfill
	\begin{figure}
		\begin{center}
            \includegraphics[width=5cm]{nnunew.jpg}
        \end{center}
	\end{figure}
	\begin{figure}
		\begin{center}
			\includegraphics[width=7cm]{nnu1.png}
        \end{center}
	\end{figure}
\end{frame}
}
\end{document}